\scanpage{17}
\mspage{15}

Les deux foncteurs composés
\[
f=k\circ *_{mn},\quad g=*_{m',n'}\circ(i\times j)
\]
dans le \deleted{\unsure{diag}} carré \unsure{pag|page} 11, dont il faut prouver qu'ils sont homotopes\unsure{,} dans $\underline{\mathrm{Hom}}(A,B)$, peuvent être \unsure{vu|vus} comme obtenus de la façon suivante, en termes de \deleted{\unsure{deux|deux tels}} foncteurs \deleted{\unsure{de}}
\diagram{\begin{tikzcd}[ampersand replacement=\&,column sep=large]
\text{(1)} \& \varphi,\psi:\Phi_{m'+n'} \arrow[r,shift left=.45ex] \arrow[r,shift right=.45ex] \& \Phi_{m+n}
\end{tikzcd}}
d'où deux foncteurs ``transposés''
\diagram{\begin{tikzcd}[ampersand replacement=\&,column sep=large]
\text{\unsure{(2)|(1)}} \& \underline{\mathrm{Hom}}(\Phi_{m,n},X) \arrow[r,"F",shift left=.5ex] \arrow[r,"G"',shift right=.5ex] \& \underline{\mathrm{Hom}}(\Phi_{m'+n'},X)
\end{tikzcd}}
lesquels foncteurs ``induisent'' $f$ et $g$. De façon précise, désignant par $a,\deleted{\unsure{c|b}}\,\unsure{b|f}$ les objets source et but de $\Phi_{m+n}$, par $c$ l'objet qui correspond au but du chemin $c$, \inserted{de longueur $m$,} à la source du chemin $c'$ de longueur $n$ qu'il s'agit de composer, enfin par $a_0,b_0$ les \unsure{élémts|éléments} source et et but de $\Phi_{m'+n'}$, les foncteurs envisagés (1) appliquent $a_0$ dans $a\unsure{{}^{\prime}|{}^{\scriptscriptstyle ?}}$, $b_0$ dans $b$, donc définissent \deleted{\unsure{chacun}} un diagramme commutatif
\diagram{\begin{tikzcd}[ampersand replacement=\&,column sep=large]
\& \Phi_{m'+n'} \arrow[r,"\varphi",shift left=.5ex] \arrow[r,"\psi"',shift right=.5ex] \& \Phi_{m+n}\\
\text{(3)} \& \underset{\text{cat. discrète}}{\unsure{\varepsilon|\mathcal E}=\{a_0,b_0\}}\arrow[u,hook,"\alpha'"]\arrow[ur,hook,"\alpha"] \&
\end{tikzcd}}
Ceci dit, la catégorie but $B$ pour $f,g$, savoir $\underline{\mathrm{Ch}}_{\Sigma,m'+n'}(X;x,z)$, s'interprète comme la sous-catégorie pleine de $\underline{\mathrm{Hom}}(\Phi_{m,n},X)$ \deleted{\unsure{dont les}} formée

% archive sheet number: 16
